\section{Finite-Length Structured SPIN}
\label{sec:finite-certificates}

The growing outer supplies the spectral control needed by the asymptotic
theorem. At a fixed length, the goal is an explicit setup failure bound,
such as $2^{-40}$, at low encoding cost.
Reusing a sampled constituent requires control of its realized spectrum,
not just its ensemble mean; see \eqref{eq:shared-constituent-enumerator}.
Concentration or selection can consume a significant part of the finite
failure budget. Randomized outers can also admit finite guarantees; the
difficulty is proof loss, not the design.

A fixed algebraic constituent removes the additional outer-spectrum failure
event. An exact spectrum is convenient, but rigorous spectrum inequalities
can serve the same interface. This section instantiates that approach
with fixed BCH constituents and the same IMT architecture. The rate-$1/2$
construction uses the same fixed inner maps.
The finite proof replaces the growing-outer argument with bounds tailored
to that selected constituent and inner.

We give certified rate-$1/2$ and rate-$1/4$ BCH constructions.
For an upper bound $U$ on distance failure, its margin is $-\log_2U$ bits.
Appendix~\ref{app:engineering} studies how outer size, message length,
state size, and step width affect the bound. Its diagnostic curves, such as
Figure~\ref{fig:finite-k-b}, guide
parameter searches; the guarantees here cover all nonzero messages.

\subsection{The Selected Construction}
\label{sec:finite-selected-construction}

We now use a binary BCH-derived constituent $\mathcal B$ of length $256$,
dimension $128$, and minimum distance at least $38$.
Its exact weight spectrum is unknown. The proof instead uses deterministic
inequalities for that spectrum, with no exceptional event for the choice
of outer.

Let $\mathcal P$ and $\mathcal Q$ be the parity extensions of the primitive binary BCH codes
of length $255$ and designed distances $37$ and $39$, respectively.
Their dimensions are $131$ and $123$. The selected constituent satisfies
$\mathcal Q\subset\mathcal B\subset\mathcal P$ and has dimension $128$.
Appendix~\ref{app:finite-bch-definition} specifies its generators and explains
why the spectrum bounds apply to the implemented code.

For $K$ message bits, choose
\[
  L:=K/128,\qquad N:=256L=2K,\qquad (t,s):=(128,19).
\]
Use $\mathcal B$ in every outer row and the structured interleaver of
Section~\ref{sec:structured-spin}. Use the IMT maps
$A:\F_2^{19}\to\F_2^{128}$ and $C:\F_2^{128}\to\F_2^{19}$ in
Table~\ref{tab:imt-maps}, also used by the asymptotic construction.
Setup $\theta$ consists of
independent uniform row permutations, independent uniform region permutations,
and independent transvections $M_i=I+u_iv_i^\top$: sample $u_i$ uniformly
from the nonzero state vectors and $v_i$ uniformly from $u_i^\perp$,
including zero.
One setup is shared by all messages.

The inner applies \eqref{eq:structured-inner} to the routed $128$-bit
groups $X_i$:
\[
  Q_0=0,\qquad Y_i=X_i+A(Q_i),\qquad
  Q_{i+1}=M_iQ_i+C(X_i).
\]
The inner and the routing distribution are therefore those of
Section~\ref{sec:structured-distance}; only the outer and the first-moment
calculation differ.

\subsection{Finite Distance Guarantee}

\begin{samepage}
\begin{theorem}[Certified BCH-256 Structured SPIN]
\label{thm:finite-bch-spin}
\leavevmode\par
For each $K\in\{2^{16},2^{18},2^{20},2^{22},2^{24}\}$, let $E_\theta$ be
the selected encoder above and put $H:=\lfloor N/10\rfloor$.
Every realization is a binary linear code of rate $1/2$, and
\begin{equation}
  \Prob_\theta\!\left[
    \exists x\ne0:\ \wt(E_\theta(x))\le H
  \right]\le U_K<2^{-40}.
  \label{eq:finite-bch-distance}
\end{equation}
The certified bounds $U_K$ have the margins reported in
Table~\ref{tab:finite-bch-margins}.
\end{theorem}
\end{samepage}

The margins range from $41.818$ to $50.189$ bits; at $K=2^{20}$
the margin is $50.062$ bits.

Thus, except with probability below $2^{-40}$ over setup, every nonzero
message has output weight greater than $N/10$.
The margin measures a certified upper bound on distance failure, not the
true failure probability or the security level of a complete application.
The theorem covers the five listed lengths, rather than every length between
them. It does not assert an asymptotic distance guarantee for fixed
$\mathcal B$ and fixed inner parameters.

\paragraph{Proof Overview.}
Let $Z_H$ count nonzero messages whose output weight is at most $H$.
Let $\mathcal T$ be the ordered row-weight classes of nonzero messages;
each class determines the number $q$ of nonzero outer rows.
Theorem~\ref{thm:class-first-moment} gives
\begin{equation}
  \Prob_\theta[Z_H>0]\le\E_\theta Z_H
  \le\sum_{\tau\in\mathcal T}
      \widehat A_\tau^{\mathrm{out}}\widehat Q_\tau(H)
  =:U_K.
  \label{eq:finite-certificate}
\end{equation}
Here $\widehat A_\tau^{\mathrm{out}}$ bounds the class multiplicity and
$\widehat Q_\tau(H)$ bounds its conditional failure probability.
No independence between different codewords is required.

The $q=1$ contribution combines a nonnegative inner transfer bound with
weighted BCH spectrum inequalities. Sparse occupations use a positive
polynomial recurrence that keeps the region-weight constraint. Dense
occupations use a comparison with independent Bernoulli inputs, together
with input and output tilts. The all-ones outer word is counted separately:
its exact multiplicity and deterministic support must not be inflated by an
ordinary weight-band envelope.

For the three largest instances, the sparse recurrence covers $2\le q\le511$,
and interval certificates cover $512\le q\le L$. The instances
$K=2^{16}$ and $K=2^{18}$ end the sparse recurrence at $q=63$ and $q=127$,
respectively, and use sharper routing-density comparisons, with separate bounds for
assignments containing high-weight rows. All contributions are summed in
exact rational arithmetic. Appendix~\ref{app:finite-bch-proof}
gives the mathematical reductions and identifies the bound inputs and
coverage checks.

\subsection{A Rate-$1/4$ Operating Point}
\label{sec:finite-quarter}

At rate $1/4$, use the fixed BCH-derived $[128,32,32]$ constituent specified
in Appendix~\ref{app:finite-quarter-definition}. Its spectrum is known
exactly. Set $K=2^{20}$, $L=K/32$, and $N=128L=4K$. Keep the same
routing distribution, expansion $A$, $(t,s)=(128,19)$, and transvection
law. Use the weight-three feedback map in
Table~\ref{tab:finite-quarter-feedback}. Only the outer code and the
feedback map change; the inner recurrence and end conditions are unchanged.

The certified targets are
$(\delta,\lambda)=(0.165,40)$ and $(0.190,30)$.

\begin{theorem}[Certified BCH-128 Structured SPIN]
\label{thm:finite-quarter-spin}
For this encoder, every setup gives a binary linear code of rate $1/4$.
For either of these two pairs $(\delta,\lambda)$,
\[
 \Prob_\theta[\exists x\ne0:\wt(E_\theta(x))\le\lfloor\delta N\rfloor]
 \le U_\delta<2^{-\lambda}.
\]
\end{theorem}

The margins are $41.048$ bits at $\delta=0.165$ and $30.033$ bits at
$\delta=0.190$. Appendix~\ref{app:finite-margin-tables} records the exact
integer cutoffs and additional digits.

The proof uses the same IMT transfer with the separately enumerated feedback
spectrum and cancellation data, followed by a complete sum over
$1\le q\le32768$. These two bounds concern different distance thresholds
for one setup distribution, not two independently selected inner designs.
